\chapter{مروری بر کارهای پیشین}
\label{chap:relatedworks}

در این فصل، به بررسی مطالعات و پژوهش‌هایی که به طور مستقیم یا غیرمستقیم با موضوع پژوهش حاضر مرتبط هستند، می‌پردازیم تا دید جامعی از آنچه تاکنون در زمینه‌های مرتبط انجام شده ارائه دهیم. ابتدا در \autoref{sec:rw_taxonomy}، دیدگاه‌های مختلف برای دسته‌بندی کارهای پیشین بیان می‌شود. در \autoref{sec:rw_criteria}، ملاک‌های ارزیابی و مقایسه تعریف می‌گردند. در \autoref{sec:rw_robust}، روش‌های مقاوم‌سازی در برابر گره‌های مخرب مرور می‌شود. در \autoref{sec:rw_hetero}، روش‌های مدیریت ناهمگونی بررسی می‌گردد. در \autoref{sec:rw_comm}، رویکردهای بهینه‌سازی ارتباطی ارائه می‌شود و در نهایت در \autoref{sec:rw_summary}، جمع‌بندی و جدول مقایسه‌ای ارائه می‌گردد.


\section{دیدگاه‌های دسته‌بندی کارهای پیشین}
\label{sec:rw_taxonomy}

پژوهش‌های مرتبط با یادگیری فدرال در محیط‌های ناهمگون را می‌توان از سه دیدگاه کلی دسته‌بندی کرد:

\subsection{دسته‌بندی بر اساس نوع ناهمگونی هدف}

اولین دیدگاه تقسیم‌بندی پژوهش‌ها بر اساس نوع ناهمگونی است که هر کار با آن مقابله می‌کند:

\begin{itemize}
    \tick \textbf{ناهمگونی داده}: پژوهش‌هایی که تمرکز اصلی آن‌ها بر مسئله توزیع \gls{NonIID} و واگرایی مدل‌های محلی است. نمونه‌های شاخص این دسته شامل \lr{FedProx} \cite{li2020federated}، \lr{SCAFFOLD} \cite{karimireddy2020scaffold} و \lr{Moon} \cite{li2021model} هستند.
    
    \tick \textbf{ناهمگونی سیستمی}: پژوهش‌هایی که به تفاوت در توان محاسباتی و شبکه‌ای گره‌ها می‌پردازند. \lr{FedProx} و \lr{Asynchronous FL} \cite{xie2019asynchronous} از این دسته‌اند.
    
    \tick \textbf{ناهمگونی بیزانتین}: پژوهش‌هایی که هدف آن‌ها مقاوم‌سازی در برابر گره‌های مخرب است. \lr{Krum} \cite{blanchard2017machine}، میانگین کوتاه‌شده \cite{yin2018byzantine} و \lr{Bulyan} \cite{el2018byzantine} از این دسته‌اند.
\end{itemize}

\subsection{دسته‌بندی بر اساس محل مداخله}

دومین دیدگاه بر این اساس است که مداخله اصلی در کجای فرایند یادگیری فدرال انجام می‌شود:

\begin{itemize}
    \tick \textbf{مداخله در تجمیع}: روش‌هایی که الگوریتم تجمیع سرور را تغییر می‌دهند (مانند \lr{Krum}، میانگین کوتاه‌شده).
    
    \tick \textbf{مداخله در آموزش محلی}: روش‌هایی که فرایند آموزش در هر گره را تغییر می‌دهند (مانند \lr{FedProx} که یک اصطلاح منظم‌ساز به تابع هدف محلی اضافه می‌کند).
    
    \tick \textbf{مداخله در ارتباطات}: روش‌هایی که به بهینه‌سازی تبادل اطلاعات بین گره‌ها و سرور می‌پردازند (فشرده‌سازی گرادیان، ارتباط ناهمزمان).
\end{itemize}

\subsection{دسته‌بندی بر اساس نیاز به اطلاعات اضافه}

سومین دیدگاه بر اساس میزان اطلاعات فراسازمانی که هر روش نیاز دارد است:

\begin{itemize}
    \tick \textbf{بدون نیاز به اطلاعات جانبی}: روش‌هایی که تنها بر اساس به‌روزرسانی‌های دریافتی در هر دور تصمیم می‌گیرند.
    
    \tick \textbf{نیاز به داده‌های اعتبارسنجی}: روش‌هایی که برای ارزیابی کیفیت به‌روزرسانی‌ها به یک مجموعه داده اعتبارسنجی در سرور نیاز دارند.
    
    \tick \textbf{نیاز به تاریخچه}: روش‌هایی که رفتار گذشته گره‌ها را نگهداری می‌کنند و در تصمیم‌گیری استفاده می‌کنند.
\end{itemize}


\section{ملاک‌های ارزیابی و مقایسه}
\label{sec:rw_criteria}

برای مقایسه منصفانه روش‌های مختلف، شش ملاک کلیدی زیر را در نظر می‌گیریم:

\begin{enumerate}
    \item \textbf{دقت نهایی مدل}: دقت مدل سراسری روی مجموعه آزمون در پایان آموزش.
    
    \item \textbf{سرعت همگرایی}: تعداد دورهای ارتباطی مورد نیاز برای رسیدن به یک دقت هدف مشخص.
    
    \item \textbf{پیچیدگی محاسباتی}: هزینه محاسباتی الگوریتم تجمیع در سرور، بر حسب تعداد گره‌ها $n$ و ابعاد مدل $d$.
    
    \item \textbf{مقاومت در برابر حمله}: توانایی حفظ دقت مدل در حضور گره‌های مخرب.
    
    \item \textbf{مقاومت در برابر ناهمگونی داده}: توانایی همگرایی در حضور توزیع \gls{NonIID}.
    
    \item \textbf{مدیریت گره‌های کُند}: توانایی ادامه آموزش بدون وابستگی به پاسخ همه گره‌ها.
\end{enumerate}


\section{روش‌های مقاوم‌سازی در برابر گره‌های مخرب}
\label{sec:rw_robust}

\subsection{\lr{Krum} و \lr{Multi-Krum}}

\lr{Krum} \cite{blanchard2017machine} اولین روش با ضمانت نظری برای یادگیری فدرال مقاوم بود. ایده اصلی آن این است که به‌روزرسانی‌های گره‌های صادق باید به یکدیگر نزدیک باشند (از نظر فاصله اقلیدسی در فضای پارامتر). بنابراین، گره‌هایی که به‌روزرسانی‌شان با سایر گره‌ها فاصله زیادی دارد، احتمالاً مخرب هستند.

تابع امتیاز \lr{Krum} برای گره $i$ به شکل زیر است:
\begin{equation}
    s(i) = \sum_{i \to j} \|\Delta_i - \Delta_j\|^2
    \label{eq:krum_full}
\end{equation}
که $i \to j$ به معنای $j$ در میان $n - f - 2$ نزدیک‌ترین همسایه $i$ است.

این روش چند محدودیت اساسی دارد. نخست، پیچیدگی محاسباتی $O(n^2 d)$ دارد که با افزایش تعداد گره‌ها یا ابعاد مدل، هزینه‌بر می‌شود. دوم، در حالت \lr{Single-Krum} (انتخاب فقط یک گره) نرخ همگرایی کُند است. سوم، در برابر حملاتی که جهت مشابه گره‌های صادق دارند (مثل حملات هم‌راستای ظریف) آسیب‌پذیر است.

\lr{Multi-Krum} \cite{blanchard2017machine} نسخه بهبودیافته است که $m > 1$ گره با کمترین امتیاز را انتخاب می‌کند. این موجب بهبود همگرایی می‌شود اما پیچیدگی محاسباتی یکسان دارد.

\subsection{میانگین کوتاه‌شده و میانه}

\lr{Yin} و همکاران \cite{yin2018byzantine} دو روش مختصه‌ای مقاوم ارائه دادند. در میانگین کوتاه‌شده، در هر مختصه، $\beta$ درصد از بزرگ‌ترین و کوچک‌ترین مقادیر حذف و میانگین مقادیر باقی‌مانده محاسبه می‌شود. پیچیدگی این روش $O(nd \log n)$ است.

در روش میانه، در هر مختصه، میانه مقادیر به‌عنوان مقدار تجمیع‌شده انتخاب می‌شود. این روش پیچیدگی $O(nd)$ دارد و از نظر مقاومت نظری قوی است، اما در عمل همگرایی کُندتری نسبت به میانگین کوتاه‌شده دارد زیرا واریانس بالاتری دارد.

\subsection{\lr{Bulyan}}

\lr{Bulyan} \cite{el2018byzantine} یک ترکیب دو مرحله‌ای است: در مرحله اول از \lr{Krum} برای انتخاب زیرمجموعه‌ای از گره‌های احتمالاً صادق استفاده می‌کند، و در مرحله دوم میانگین کوتاه‌شده را روی این زیرمجموعه اعمال می‌کند. این ترکیب از نظر نظری قوی‌تر است، اما پیچیدگی محاسباتی آن نیز بالاتر است و نیاز به $n \geq 4f + 3$ دارد.

\subsection{\lr{FLTrust}}

\lr{FLTrust} \cite{cao2020fltrust} یک رویکرد متفاوت را دنبال می‌کند. سرور یک مجموعه داده کوچک و قابل‌اعتماد دارد و مدل خود را روی آن آموزش می‌دهد. به‌روزرسانی هر گره بر اساس شباهت کسینوسی آن با به‌روزرسانی سرور وزن‌دهی می‌شود. این روش در عمل نتایج خوبی داشته، اما نیاز به داده اعتبارسنجی در سرور دارد که در بسیاری از سناریوهای واقعی در دسترس نیست.

\subsection{\lr{FedDF} و \lr{FedBE}}

این دو روش \cite{lin2020ensemble, chen2021fedbe} از تقطیر دانش (\lr{Knowledge Distillation}) برای تجمیع استفاده می‌کنند. به جای میانگین مستقیم پارامترهای مدل، مدل سراسری را با استفاده از خروجی مدل‌های محلی روی یک داده کوچک در سرور آموزش می‌دهند. این رویکرد در برابر ناهمگونی داده قوی است اما نیاز به داده در سرور دارد.


\section{روش‌های مدیریت ناهمگونی}
\label{sec:rw_hetero}

\subsection{\lr{FedProx}}

\lr{FedProx} \cite{li2020federated} با افزودن یک اصطلاح منظم‌ساز به تابع هدف محلی، واگرایی مدل‌های محلی از مدل سراسری را محدود می‌کند:

\begin{equation}
    \min_w F_i(w) + \frac{\mu}{2}\|w - w^t\|^2
    \label{eq:fedprox}
\end{equation}

پارامتر $\mu$ شدت منظم‌سازی را کنترل می‌کند. این روش همچنین امکان آموزش ناقص در گره‌های کُند را فراهم می‌کند که برای محیط‌های ناهمگون مفید است. \lr{FedProx} در برابر ناهمگونی داده بهتر از \lr{FedAvg} عمل می‌کند، اما مکانیزمی برای مقاوم‌سازی در برابر گره‌های مخرب ندارد.

\subsection{\lr{SCAFFOLD}}

\lr{SCAFFOLD} \cite{karimireddy2020scaffold} ریشه واگرایی مدل‌های محلی در توزیع \gls{NonIID} را در نوسانات واریانس گرادیان می‌داند. این روش از متغیرهای کنترلی $c_i$ و $c$ برای اصلاح به‌روزرسانی‌های محلی استفاده می‌کند:

\begin{equation}
    w_i \leftarrow w_i - \eta (g_i - c_i + c)
    \label{eq:scaffold}
\end{equation}

این روش همگرایی نظری قوی‌تری دارد اما دو برابر \lr{FedAvg} ارتباط نیاز دارد.

\subsection{\lr{Moon}}

\lr{Moon} \cite{li2021model} از یادگیری مقایسه‌ای (\lr{Contrastive Learning}) برای بهبود آموزش محلی استفاده می‌کند. ایده اصلی این است که نمایش مدل محلی باید به نمایش مدل سراسری نزدیک باشد و از نمایش مدل محلی قبلی فاصله داشته باشد.

\subsection{\lr{FedNova}}

\lr{FedNova} \cite{wang2020tackling} مشکل واگرایی را از منظر نرخ تغییر مؤثر مدل مورد بررسی قرار می‌دهد. این روش به‌روزرسانی‌های گره‌ها را بر اساس تعداد واقعی مراحل \gls{SGD} که هر گره انجام داده نرمال‌سازی می‌کند.


\section{رویکردهای بهینه‌سازی ارتباطی}
\label{sec:rw_comm}

\subsection{ارتباط ناهمزمان}

در یادگیری فدرال ناهمزمان \cite{xie2019asynchronous}، سرور منتظر پاسخ همه گره‌های انتخاب‌شده نمی‌ماند و به محض دریافت به‌روزرسانی از هر گره، مدل را به‌روز می‌کند. این رویکرد مشکل گره‌های کُند را حل می‌کند، اما ممکن است به دلیل استفاده از اطلاعات قدیمی (به‌اصطلاح \lr{Staleness})، پایداری آموزش را تحت تأثیر قرار دهد.

\subsection{فشرده‌سازی گرادیان}

یکی از راه‌های کاهش حجم ارتباطات در یادگیری فدرال، فشرده‌سازی به‌روزرسانی‌های ارسالی است. روش‌های متداول شامل کوانتیزاسیون \cite{reisizadeh2020fedpaq} و هرس (\lr{Pruning}/\lr{Sparsification}) \cite{lin2017deep} هستند. این روش‌ها ابعاد داده ارسالی را کاهش می‌دهند اما ممکن است خطا وارد کنند.

\subsection{انتخاب هوشمند گره}

برخی از پژوهش‌ها به جای انتخاب تصادفی گره‌ها در هر دور، از معیارهای هوشمندانه‌تری استفاده می‌کنند \cite{lai2021oort}. معیارهایی نظیر حجم داده محلی، توان محاسباتی، یا تنوع داده‌ای می‌توانند در انتخاب گره‌های مشارکت‌کننده مؤثر باشند.


\section{جمع‌بندی و مقایسه}
\label{sec:rw_summary}

بررسی کارهای پیشین نشان می‌دهد که هیچ روش موجودی به‌طور هم‌زمان سه چالش اصلی محیط‌های ابری چندگانه را حل نمی‌کند: مقاومت در برابر گره‌های مخرب، مدیریت ناهمگونی داده، و مدیریت گره‌های کُند. به‌علاوه، روش‌هایی که در برابر گره‌های مخرب قوی هستند (مانند \lr{Krum}) پیچیدگی محاسباتی بالایی دارند که در محیط‌های ابری با هزینه مستقیم همراه است.

روش پیشنهادی این پژوهش، \gls{TWAC}، با استفاده از امتیازات اعتماد تاریخی و یک مکانیزم دو-مؤلفه‌ای سبک‌وزن، به همه این چالش‌ها به‌صورت یکپارچه پاسخ می‌دهد. جدول \ref{tab:related_works_comparison} مقایسه جامع روش‌های مختلف را نشان می‌دهد.

\begin{table}[t]
    \centering
    \renewcommand{\arraystretch}{1.8}
    \caption{مقایسه روش‌های مختلف یادگیری فدرال بر اساس ملاک‌های کلیدی}
    \label{tab:related_works_comparison}
    \begin{tabular}{lcccccc}
        \toprule
        \textbf{روش} & 
        \textbf{پیچیدگی} & 
        \textbf{مقاومت} & 
        \textbf{ناهمگونی} & 
        \textbf{گره کُند} & 
        \textbf{داده} & 
        \textbf{حافظه} \\
        & & \textbf{بیزانتین} & \textbf{داده} & & \textbf{سرور} & \textbf{تاریخی} \\
        \midrule
        \lr{FedAvg}         & $O(nd)$        & \ding{55} & متوسط     & \ding{55} & \ding{55} & \ding{55} \\
        \lr{FedProx}        & $O(nd)$        & \ding{55} & خوب       & بله       & \ding{55} & \ding{55} \\
        \lr{SCAFFOLD}       & $O(nd)$        & \ding{55} & خوب       & \ding{55} & \ding{55} & \ding{55} \\
        \lr{Krum}           & $O(n^2 d)$     & قوی       & ضعیف      & \ding{55} & \ding{55} & \ding{55} \\
        میانگین کوتاه‌شده  & $O(nd\log n)$  & خوب       & متوسط     & \ding{55} & \ding{55} & \ding{55} \\
        میانه               & $O(nd)$        & خوب       & متوسط     & \ding{55} & \ding{55} & \ding{55} \\
        \lr{Bulyan}         & $O(n^2 d)$     & قوی       & ضعیف      & \ding{55} & \ding{55} & \ding{55} \\
        \lr{FLTrust}        & $O(nd)$        & قوی       & متوسط     & \ding{55} & بله       & \ding{55} \\
        \lr{FedNova}        & $O(nd)$        & \ding{55} & خوب       & بله       & \ding{55} & \ding{55} \\
        \lr{TWAC} (این پژوهش) & $O(nd)$    & قوی       & خوب       & بله       & \ding{55} & بله       \\
        \bottomrule
    \end{tabular}
\end{table}

جدول \ref{tab:related_works_comparison} به روشنی نشان می‌دهد که روش \gls{TWAC} تنها روشی است که بدون نیاز به داده اعتبارسنجی در سرور، مقاومت در برابر گره‌های مخرب، مدیریت ناهمگونی داده و مدیریت گره‌های کُند را با هم ترکیب می‌کند، و این کار را با پیچیدگی محاسباتی بهینه $O(nd)$ انجام می‌دهد.
